\section{Standardisation et transition bibliographique}
La transition bibliographique a pour objectif de mieux faire remonter les données des bibliothèques dans les moteurs de recherches, par des moyens différents. La philosophie des notices bibliographies ont changé comme l'explique Marianne Clatin. \enquote{On ne catalogue plus pour que l'usager vienne sur notre catalogue; on transforme nos notices en données pour qu'elles se trouvent directement sur le chemin de l'usager, c'est-à-dire dans les moteurs de recherche généralistes (Google, etc.) grâce au référencement sémantique(SEO)}.

L'enjeu actuel pour les institutions patrimoniales, dont la bibliothèque Cujas, est d'avoir un meilleur référencement, c'est-à-dire une position parmi les premières pages retournées par Google à la suite d'une recherche. Aujourd'hui, le web a pour effet immédiat de créer une économie de l'abondance de l'information. La bibliothèque n'est donc plus un passage obligé pour accéder aux documents. Toute démarche orientée utilisateur doit donc prendre en compte comme paramètre le besoin de visibilité, permettant de capter l'attention des internautes.

\subsubsection{L'application d'IFLA-LRM aux objets complexes}
Le passage de la logique de \enquote{notices isolées} à celle de \enquote{données} interconnectées pour s'adapter au parcours de l'usager, lui permettant de faire des recherches par entités (personnes, lieux, oeuvres). Cette déconstruction de la \enquote{notice plate} au profit du modèle OEMI (Oeuvre, Expression, Manifestation et Item), décrivant et structurant les ressources bibliographiques lors de cette transition bibliographique. 

Lorsque l'on parle d'objets complexes, on désigne les périodiques représentant un objet difficile à définir dans son ensemble pour les institutions patrimoniales. Il ne s'agit pas uniquement du logiciel utilisé par les bibliothèques pour les exposer. La bibliothèque Cujas, à l'origine, a choisi de développer sa première bibliothèque numérique sous XTF avec la refonte les experts métiers envisagent une bibliothèque qui acceptent différentes types de documents. 

L'attente interne est d'intégrer les périodiques dans sa future bibliothèque. Il faut que une structure capable de \enquote{modéliser finement l'arborescence multi-niveaux propre aux revues scientifiques (Titre de la revue > Année > Fascicule/Volume > Article)}, avec une indexation détaillée poussée jusqu'au niveau de l'article. L'application des entités du modèle conceptuel IFLA-LRM permettrait de traduire cette complexité éditoriale et de combler cette lacune présente de CujasNum. Aujourd'hui, le document numérisé n'est plus un simple objet physique, mais il possède une architecture de données manipulable. Comme l'a analysé Jean-Michel Salaün, \enquote{l'utilisation de référentiels et de données modélisées en graphe pour créer du lien entre des ressources d'origine différentes, dans différents formats, qui traitent des mêmes sujets}. Ce nouveau modèle OEMI se présente avec la promesse de transformer en données interopérables tout l'information manipulée par les institutions patrimoniales. 

Les liens entre les informations se feraient à l'aide de triplet permettant à la machine de différencier les éléments grâce à l'URI, favorisant ainsi les liens de plusieurs manières des informations portant sur un même document. Il devint possible de lier une thèse de droit, par exemple, à un exemplaire physique, à un substitut numérique, à la personne qui l'a écrite, celle qui l'a dirigée en définissant précisément les relations de chacun de ces éléments. Cela permettrait de répondre à une des attentes de la direction de la BIU Cujas qui est de vouloir intégrer des réseaux nationaux en traitant de la donnée massivement. Ces liens permettent de décrire le plus précisément possible une chose (document, contenu, personne,etc.), permettantle \enquote{passeport du document}, comme l'expliquait Olesea en signalant plus il y a d'informations mieux c'est pour l'usager. L'une des solutions est donc de constituer des notices d'autorité.

\subsubsection{L'alignement d'autorités et l'URI pivot}
La majorité des notices bibliographiques sont encodés sous Dublin Core simple, qui les appauvrissent car uniquement cinq balises sur les quinze possible sont généralement utilisés. Il est important de les enrichir, notamment celles sur les livres anciens. Selon le constat d'Olesea Dubois, la richesse des informations est nécessaire pour aider l'usager à comprendre le document qu'il voit numériser aussi bien le visuel que sa matérialité. Le document numérisé est uniquement ce que voit l'usager sur la bibliothèque numérique. En effet, \enquote{la bibliothèque numérique est la seule chose que l'utilisateur voit, c'est la vitrine du document, plus on est bavard dans les descriptions}, meilleure est la compréhension intégrale du document. 

La transition d'une logique de notices isolées vers un réseau sémantique repose sur un changement dans l'art de décrire les ressources. Dans l'écosystème du Web de données, l'autorité n'est plus une chaîne de caractère figée dans un catalogue local, mais devenu une entité dynamique identifiée par une URI. Pour modéliser cette interconnexion, Emmanuelle Bermès, avec la collaboration d'Antoine Isaac et Gautier Poupeau, mobilise la métaphore du \enquote{hub and spoke}. Selon les auteurs, les référentiels, thésaurus et listes d'autorités contrôlées agissent comme le point nodal permettant de créer un ou plusieurs points de contact entre les données. \enquote{Les référentiels ou vocabulaires contrôlés qui fournissent les réservoirs d’URI réutilisables sont donc appelés à jouer un rôle essentiel dans le Web de données. Ils fournissent en effet les points d’ancrage qui vont transformer les graphes isolés de chaque institution en un graphe global, relié par des points de contacts communs}. Ces points de contact permettent de naviguer sans contrainte entre les données, en utilisant les URI. La particularité de cette métaphore est qu'il n'est pas nécessaire d'exprimer les données suivant le même modèle pour qu'elles puissent communiquer entre elles.

Appliquée à la refonte de la bibliothèque numérique de la bibliothèque Cujas, cette conceptualisation constitue un ancrage méthodologique fondamental. Elle démontre qu'il est inutile de chercher à unifier l'intégralité des structures de leurs bases de données locales ou de développer des formats propriétaires complpexes. Il suffit d'aligner les notices bibliographiques ou les fiches biographiques sur les référentiels nationaux de l'Enseignement Supérieur et de la Recherche (ESR)n principalement IdRef pour les personnes et RAMEAU pour les sujets. Sur le plan technique, cet alignement résout deux obstacles majeurs de l'accès à l'information : \enquote{l'homonymie (en distinguant deux auteurs aux patronymes identiques grâce à leur identifiant IdRef unique)}ou \enquote{la synonymie (en regroupant sous une même entité conceptuelle des termes différents ou des varitations linguistiques)}. Pour faire cet alignement de masse, deux outils existent Bibliostratus  et OpenRefine, avec des objectifs différents. 

L'alignement massif est une condition \textit{sine qua none} pour lier de manière automatisée et à grandes échelles les données locales aux grands référentiels du Web (IdRef, VIAF ou Wikidata), en brisant les silos du Web, en enrichissant les catalogues des bibliothèques sans l'intervention humaine. Un autre avantage de l'alignement est qu'il améliore l'expérience utilisateur permettant de créer des recherche par rebonds, des recherches multilinguistes et d'améliorer la visibilité des bibliothèques. Ces trois conditions sont des attentes fortes de la bibliothèque Cujas à l'issue de cette refonte qu'elle pourra proposer à ses utilisateurs. Un autre enjeu pour l'établissement est de parvenir à trouver un mode d'indexation plus adapté afin de constituer des facettes plus pertinentes pour eux également.

\subsubsection{L'alignement entre l'indexation Hauville et RAMEAU}
Les exigences de refonte sémantique formulées par le personnel de Cujas touchent également à la rationalisation de l'indexation thématique, aujourd'hui sujette à un bruit documentaire important. L'ancien responsable avait préconisé une \enquote{limitation stricte stricte à cinq indexations maximum par ouvrage}, afin de concentrer la prolifération lexicale et d'améliorer la \enquote{trouvabilité} des documents pour les chercheurs. Cette volonté de rationalisation, qui invite à se baser sur le ressenti de l'usager plutôt que d'empiler des concepts techniques complexes déconnectés des recherches réelles, rejoint la problématique de l'interopérabilité. Entre la standardisation nationale portée par le référentiel RAMEAU et la finesse de l'indexation Hauville, la négociation d'un alignement sémantique en SKOS s'impose comme une solution majeure pour répondre à ses enjeux de modernisation. Grâce à l'alignement sémantique, la bibliothèque Cujas permettrait de faire évoluer le catalogue sans tout réindexer manuelle à chaque fois qu'un terme de droit change. 

L'alignement sémantique en SKOS permettrait à la bibliothèque de conserver l'indexation Hauville, comme un thésaurus autonome et spécialisé en droit, et créer un graphe d'alignement SKOS indépendant permettant de faire le pont entre les deux. Cela réduirait alors le bruit observé lorsqu'un étudiant effectue une recherche, par exemple, dans CujasNum en cherchant à utiliser les facettes à sa disposition. 

Le choix de la bibliothèque Cujas, face à la création de sa propre indexation, résultait du constat que le référentiel RAMEAU n'est pas toujours entièrement adapté, manquant à certain moment de précision. La réfonte ne doit pas entraîner un écrasement ou l'abandon de l'indexation Hauville, mais résulte d'un travail scientifique de longue haleine. \enquote{L'harmonisation impose des choix qui peuvent entrer en conflit avec les besoins documentaires spécifiques, et pose la question de trouver comment accorder les richesses des connaissances de chaque institution concernée avec la nécessité d'avoir un minimum d'uniformité lorsqu'il s'agit de coopéreer à grande échelle}. En effet, la création d'un alignement sémantique permettrait de faire un travail de mapping entre l'indexation Hauville et le référentiel RAMEAU. En effet, comme l'a analysé Maëlys Gioan dans son étude, en expliquant qu'aligner, \enquote{c'est relier explicitement des concepts équivalents ou proches, structurer leurs correspondancess à l'aide des proprités (en SKOS, skos:axactMatch, closeMatch, broadMatch ou related)}. Cette pluralité des langages n'est pas un problèmes informatique mais une réalité humaine et scientifique, répondant à des valeurs ou des missions différentes.

L'indexation s'impose comme une brique incontournable de la constitution d'une bibliothèque numérique. Elle a besoin d'avoir fait l'objet d'un travail scientifique intellectuel conséquent car à partir de ce travai on crée les facettes. Ces facettes ont un \enquote{rôle cognitif}. Comme l'a étudié Elina Leblanc, elle explique que \enquote{les facettes reflètent la manière dont la bibliothèque numérique a pensé ses collections et aident ainsi l'utilisateur à apprendre empiriquement l'organisation des collections et à naviguer parmi les contenus}. Une indexation trop complexe ou mal hiérarchisée, comme le \enquote{problème des tirets}, perturbe cet apprentissage sémantique. Une bonne indexation est nécessaire, afin d'éviter la surabondance informatique. Rationaliser le nombre d'entrées à cinq pour l'indexation est nécessaire pour éliminer le \enquote{bruit}, réduisant la surchange informationnelle auprès des usagers.
\bigskip

En définitive, la normalisation sémantique et l'adoption de protocoles standards permet à la bibliothèque Cujas de s'assurer que ses collections patrimoniales ne sont plus de simples fichiers isolés, mais des entités interconnectées, circulant librement sur le Web de données. Pour autant, cette rigueur informatique et sémantique ne saurait constituer une fin en soi. Si le Web sémantique est conçu pour être lu et interprété par des machines, la bibliothèque numérique, elle, est crée pour des êtres humains. Les graphes de triplets (RDF) et les serveurs IIIF ne prenneent leur véritable valeur sicnetifique que lorsque cela se traduit en bénéfices concrets pour l'usager final. Toutefois, lorsque l'on conçoit une interface idéale sur le plan technique mais qui fassent abstractions des réalités cognitives, des compétences informationnelles et de la façon dont les équipes travaillent. Il est important que la future plateforme de la bibliothèque Cujas replacea l'usager au centre de la réflexion en analysant leurs attentes.